\section{Support and strip arguments with predecessor loss}
\label{ms:legacy-support:section}

The finite-index proofs use support information before they enter a
Fourier or conditional-count calculation. Some of those implications
allow a positive predecessor loss, a positive score excess, or an index
below $42$. They are not consequences of the zero-loss theorem in
Section~\ref{ms:support}. This appendix supplies their analytic arguments
with their literal index domains. The numerical polynomial inequalities,
moment rows and threshold enclosures remain separately stated certificate
premises.

We use the attained positive closed pair and contact identity of
Lemma~\ref{ms:contact}. Thus the law is fixed throughout every deletion,
moment estimate and support contraction; $H$ is its raw sum threshold.
The structural bounds \eqref{ms:structural-classical}, the sign-moment
inequality of Lemma~\ref{ms:sign-moment}, and, when used to acquire a
signed moment row, Theorem~\ref{ms:sharp-third}, retain their stated
hypotheses. Write $C=C_n$ for the attained score, $\rho=\E|X|^3$ and
$\beta=\E[X|X|]$. Every certificate below must contain the same law's
parameters and the whole closed domain on which it is applied.

\subsection{The scale correction}
Write $\phi_0=\phi(0)$.
We use the rational comparisons of Lemma~\ref{ms:gaussian-constants}.
For $n>1$ put $s_n=\sqrt{n/(n-1)}$ and
$e_n=n(s_n-1)-1/2$.
\begin{lemma}\label{ms:legacy-support:scale}
For every integer $n\ge36$, one has $0<e_n<0.384/n$.
\end{lemma}
\begin{proof}
Put $x=1/n$ and $c=0.384=48/125$. Direct expansion gives
\[
 (1-x)(1+x/2+cx^2)^2-1
 =x^2\{(2c-3/4)-(c+1/4)x+(c^2-c)x^2-c^2x^3\}.
\]
The bracket is decreasing for $x\ge0$, since $0<c<1$.
At $x=1/36$ its exact rational value is $823/4050000>0$.
Squaring positive quantities gives
$s_n<1+x/2+cx^2$ for $0<x\le1/36$. Similarly
$(1-x)(1+x/2)^2-1=-3x^2/4-x^3/4<0$, which gives the lower bound.
\end{proof}

\subsection{A signed Gaussian envelope}
Define, for every real $H,t$ and $n>1$,
\[
 \mathcal H_n(H,t)=n^{3/2}\{\Phi((H-t)/\sqrt{n-1})-
             \Phi(H/\sqrt n)\}
       +nt\phi(H/\sqrt n)+\tfrac H2(t^2-1)\phi(H/\sqrt n).
\]
\begin{lemma}[Direct support at every index]\label{ms:legacy-support:direct}
Let \(n\ge2\), and let \((P,H)\) be an unrestricted attained
positive closed maximizing pair, with \(C=C_n>0\).
Put \(s=\sqrt{n/(n-1)}\), \(e=n(s-1)-1/2\),
\(A=\sqrt{n-1}\,\Delta_{n-1}(P)/\rho\), and \(d_P=n(A-C)\).
Suppose \(q\ge0\), \(S>0\), \(E\ge0\), \(r>0\), \(b\ge0\),
\((d_P)_+\le q\), \(s\le S\), \(e\le E\),
\(\rho\le r\), \(|\beta|\le b\), and \(0<c\le C\).
Define
\[
 Q(v)=-v^3+\tfrac32rv^2+3bv+rE,\qquad
 L(v)=\tfrac25(v^2/2+1/2+E)+cQ'(v).
\]
If \(R>0\) satisfies
\[
 Q(R),Q'(R),Q''(R),L(R),L'(R)<0,\qquad 1/5-3c<0,
\]
and, for both \(\eta\in\{-1,1\}\),
\[
 \sup_{H\in\R}\mathcal H_n(H,\eta R)+cQ(R)+rSq<0,
\]
then \(\supp P\subset(-R,R)\).
\end{lemma}
\begin{proof}
The contact identity~\eqref{ms:contact-formula}, the active
discrepancy \(C\rho/\sqrt n\), and the predecessor estimate
for this same law imply, at every support point \(t\),
\[
 0\le\mathcal H_n(H,t)
 +C\{-|t|^3+\tfrac32\rho t^2+3\beta t+\rho e\}
 +\rho s d_P.
\]
Here \(n(sA-C)=s d_P+n(s-1)C\); the scale correction has
no omitted term. Also \(A\le C_{n-1}\), hence
\((d_P)_+\le n(C_{n-1}-C)_+\).
Put \(z=H/\sqrt n\), \(y=(H-t)/\sqrt n\).
Exact differentiation gives
\[
 \partial_t\mathcal H_n=n\phi(z)+Ht\phi(z)-ns\phi(sy).
\]
Taylor's theorem and \(\|\phi''\|_\infty=\phi(0)\) give
\( |n\phi(y)-n\phi(z)-Ht\phi(z)|\le\phi(0)t^2/2\).
Moreover \(\partial_s[s\phi(sy)]=-\phi''(sy)\), so
\(|s\phi(sy)-\phi(y)|\le\phi(0)(s-1)\). Therefore
\[
 |\partial_t\mathcal H_n(H,t)|
 \le\phi(0)(t^2/2+1/2+e)
 \qquad(n\ge2,\ H,t\in\R).
\]
The norm identity follows by maximizing
\((x^2-1)e^{-x^2/2}\); its interior extrema occur at
\(0,\pm\sqrt3\).
Since \(Q'''=-6\), its first three gates imply \(Q(v)<0\)
for every \(v\ge R\). At \(t=\eta v\), the bracket multiplying
\(C\) is at most \(Q(v)\), so its multiplier may be replaced
by \(c\). The signed predecessor term is at most \(rSq\).
The contact expression is consequently bounded by
\(\mathcal H_n(H,\eta v)+cQ(v)+rSq\), whose derivative
in \(v\) is at most \(L(v)\). Since
\(L''=2/5-6c<0\) and \(L(R),L'(R)<0\), this envelope strictly
decreases on the entire closed ray. Its strict endpoint upper
bound contradicts contact on both rays, including endpoints,
exterior atoms and exterior accumulation points.
\end{proof}

\paragraph*{The whole threshold line.}
No sampled threshold maximum is needed in this lemma.
For \(R>0\), Gaussian differentiation gives
\[
 |\partial_H^2\mathcal H_n(H,R)|\le M,\qquad
 M=\frac{\sqrt n}{4}\left(\frac n{n-1}+1\right)
       +\frac25R+\frac{2|R^2-1|}{5\sqrt n}.
\]
This follows from \(\|\phi'\|_\infty<1/4\),
\(\|\phi''\|_\infty<2/5\), and
\(\|\phi'''\|_\infty\le2/\pi<4/5\).
For a closed threshold cell of midpoint \(m\) and halfwidth \(w\),
Taylor's theorem bounds both signs by
\[
 \pm\mathcal H_n(H,R)
 \le\pm\mathcal H_n(m,R)
       +|\partial_H\mathcal H_n(m,R)|w+Mw^2/2.
\]
If \(H_0\ge5\sqrt n\), \(H_0-R\ge5\sqrt{n-1}\), and
\(\tau\ge\Phi(-5)\), Mills' bound and monotonicity of
\(H\phi(H/\sqrt n)\) give, on both threshold tails,
\[
 |\mathcal H_n(H,R)|\le n\sqrt n\,\tau
 +(nR+H_0|R^2-1|/2)\phi(H_0/\sqrt n),\qquad |H|\ge H_0.
\]
Finally \(\mathcal H_n(H,-R)=-\mathcal H_n(-H,R)\).
Finite rational CDF, density and square-root enclosures on all
closed cells, together with these two tails, therefore certify
the displayed whole-line endpoint bounds.

In particular the positive-budget support certificate at index 27
uses
\[
 \begin{gathered}
 R=13/5,\quad q=813/2000,\quad c=8171/18000,\quad
 r=6413121/5478500,\\
 b=326352124353/10^{12},\quad E=14331929737/10^{12},\quad
 S=1019049330731/10^{12}.
 \end{gathered}
\]
The contrary hypothesis \(C_{27}\ge \CE+2/45\), the finite
cap and the signed-moment inequality supply these bounds.
Its retained numerical premise checks every polynomial gate and both strict endpoint inequalities
on 580 adjacent closed threshold cells covering \([-29,29]\),
and on the two tails. Conditional on those complete checks and the stated moment bounds, this support
implication applies at 27 without extending an estimate proved only for \(n\ge36\).

Write $v=n-1$, $s=\sqrt{n/v}$, $u=|t|$, and
\[
 A=\frac{u^3-3u}{6},\quad B=\frac{u^4-6u^2+3}{24\sqrt v},
 \quad D=\frac{u^5-15u}{120v}.
\]
\begin{lemma}[Fifth-order identity with a uniform finite remainder]
\label{ms:legacy-support:fifth}
If $u>\sqrt3$, then, with $z_v=H/\sqrt v$,
\[
 \frac{\mathcal H_n(H,\pm u)}{s^3}
 =\pm A(1-z_v^2)\phi(z_v)+B\phi^{(3)}(z_v)
                  \mp D\phi^{(4)}(z_v)+R,
\]
where the sign of the remainder need not be fixed and
\[
 |R|\le E^{(5)}_v(u):=
 \frac{51u^6}{14400v^{3/2}}
 +\frac{51(3u^2-2)}{960v^{3/2}}+\frac{u}{8v^2}.
\]
This statement is uniform over every real $H$.
\end{lemma}
\begin{proof}
Let $F_w(x)=\Phi(x/\sqrt w)$, with primes denoting $x$ derivatives.
The heat identity $\partial_w F_w=F_w''/2$ and ordinary Taylor's theorem
give, with unlabelled derivatives at $(v,H)$,
\begin{align*}
 F_{v+1}&=F_v+F_v''/2+F_v^{(4)}/8+R_0,
 &|R_0|&\le\sup_{w\in[v,v+1]}\|F_w^{(6)}\|_\infty/48,\\
 F_{v+1}'&=F_v'+F_v'''/2+F_v^{(5)}/8+R_1,
 &|R_1|&\le\sup_{w\in[v,v+1]}\|F_w^{(7)}\|_\infty/48,\\
 F_{v+1}''&=F_v''+F_v^{(4)}/2+R_2,
 &|R_2|&\le\sup_{w\in[v,v+1]}\|F_w^{(6)}\|_\infty/8.
\end{align*}
Expand $F_v(H-t)$ through degree five, with remainder at most
$u^6\|F_v^{(6)}\|_\infty/720$. In
\[
 \frac{\mathcal H_n}{n^{3/2}}=
 F_v(H-t)-F_{v+1}(H)+tF_{v+1}'(H)
             -\frac{t^2-1}{2}F_{v+1}''(H),
\]
the first three derivative coefficients cancel. The next three are
\[
 \left(\frac t2-\frac{t^3}6\right)F_v'''
 +\frac{t^4-6t^2+3}{24}F_v^{(4)}
 +\left(\frac t8-\frac{t^5}{120}\right)F_v^{(5)}.
\]
The remaining absolute value is bounded by
\[
 \left\{\frac{u^6}{720}+\frac1{48}+\frac{u^2-1}{16}\right\}
       \sup_{w\in[v,v+1]}\|F_w^{(6)}\|_\infty
 +\frac u{48}\sup_{w\in[v,v+1]}\|F_w^{(7)}\|_\infty.
\]
Fourier inversion of the Gaussian, differentiated under its absolutely
integrable Fourier transform, gives
\[
 \|\phi^{(k)}\|_\infty
 \le\frac1\pi\int_0^\infty y^k e^{-y^2/2}\,dy.
\]
Integration by parts yields
$\|\phi^{(5)}\|_\infty\le8/\pi<2.55$ and
$\|\phi^{(6)}\|_\infty\le15\phi_0<6$.
Since $F_w^{(j)}=w^{-j/2}\phi^{(j-1)}$, the displayed remainder
becomes $s^3E^{(5)}_v(u)$ after multiplication by $n^{3/2}$.
The three retained derivative coefficients give the identity.
\end{proof}

\begin{lemma}[A lower tangent bound for the fourth Gaussian derivative]
\label{ms:legacy-support:gaussian-tangent}
For every real $z$,
\[
 \phi^{(4)}(z)\ge\phi_0(3-15z^2/2).
\]
\end{lemma}
\begin{proof}
The same Gaussian Fourier integral gives
$\phi^{(4)}(z)=\pi^{-1}\int_0^\infty
y^4e^{-y^2/2}\cos(yz)\,dy$.
Use the globally valid inequality $\cos w\ge1-w^2/2$ inside
this integral. Its fourth and sixth Gaussian moments are $3\phi_0$
and $15\phi_0$, respectively.
\end{proof}

\begin{lemma}[Uniform signed fifth-order envelope]\label{ms:legacy-support:envelope}
Suppose $u>\sqrt3$ and $A,B,D\ge0$. If
$W=A-15D/2>0$, put
\[
 T=A-3D+\frac{9B^2}{4W}.
\]
If $T\ge A/2+2B+3D$, then, for every real $H$ and either sign,
\[
 \mathcal H_n(H,\pm u)\le s^3\{\phi_0 T+E^{(5)}_v(u)\}.
\]
\end{lemma}
\begin{proof}
Put $z_v=H/\sqrt v$. For positive $t=u$ and $|z_v|\le1$, divide the retained terms by
$\phi_0$ and use the preceding tangent bound to obtain
\[
 A-3D-(A-15D/2)z_v^2+3B|z_v|
 \le A-3D+9B^2/(4W).
\]
Here $(1-z_v^2)\phi(z_v)\le(1-z_v^2)\phi_0$ and
$|\phi^{(3)}(z_v)|\le3|z_v|\phi_0$ on this interval.
For $|z_v|\ge1$ the first term is nonpositive, and the Fourier bounds
$\|\phi^{(3)}\|_\infty<2\phi_0$ and
$\|\phi^{(4)}\|_\infty\le3\phi_0$ give $2B+3D$.
For negative $t=-u$, the first term has global upper bound
$2e^{-3/2}A\phi_0<A\phi_0/2$, and the other terms give
$(2B+3D)\phi_0$. The hypothesis on $T$ covers both exterior cases.
The inequality $2e^{-3/2}<1/2$ follows from four positive terms of
the exponential series at $3/2$. Finally add the uniform remainder.
\end{proof}

\subsection{A direct Gaussian bound on both complete rays}
\label{ms:legacy-support:far-gaussian-section}

\begin{lemma}[A global cubic Gaussian bound]\label{ms:legacy-support:far-gaussian}
For every integer $n\ge2$ and all real $H,t$,
\[
 |\mathcal H_n(H,t)|<\frac{|t|^3}{15}+\frac{|t|}{5}+\frac12.
\]
\end{lemma}
\begin{proof}
Put $z=H/\sqrt n$, $u=z-t/\sqrt n$, and
$f(\lambda)=\Phi(u/\sqrt{1-\lambda})$. Taylor expansion in the spatial
variable through degree two gives
\[
 n^{3/2}\{\Phi(u)-\Phi(z)\}
 =-nt\phi(z)-\tfrac12Ht^2\phi(z)+R_{\rm tr},\qquad
 |R_{\rm tr}|\le\phi(0)|t|^3/6,
\]
using $\|\phi''\|_\infty=\phi(0)$. Expansion of $f$ through degree
one at zero gives
\[
 n^{3/2}\{f(1/n)-f(0)\}
 =\tfrac12\sqrt n\,u\phi(u)+R_{\rm sc}.
\]
For $0\le\lambda\le1/2$, writing
$r=u/\sqrt{1-\lambda}$, differentiation gives
\[
 f''(\lambda)=\frac{r(3-r^2)\phi(r)}{4(1-\lambda)^2}.
\]
The exact maxima of $|r|\phi(r)$ and $|r|^3\phi(r)$ are
$\phi(0)e^{-1/2}<1/4$ and
$3\sqrt3\,\phi(0)e^{-3/2}<1/2$ respectively. These strict rational
bounds follow from $\phi(0)<2/5$, $e>8/3>64/25$,
$e^{3/2}>21/5$, and $3\sqrt3<21/4$; the exponential inequalities follow
from finite positive terms of their series. Therefore
$\|f''\|_{[0,1/2]}<5/4$, and
\[
 |R_{\rm sc}|<\frac5{8\sqrt n}\le\frac5{8\sqrt2}<\frac12.
\]
Substitution in the definition of $\mathcal H_n$ cancels the terms in
$nt$ and $Ht^2$, leaving
\[
 \mathcal H_n(H,t)=R_{\rm tr}+R_{\rm sc}
       +\tfrac12\sqrt n\{u\phi(u)-z\phi(z)\}.
\]
Since $(x\phi(x))'=-\phi''(x)$, the last term has absolute value at
most $\phi(0)|t|/2<|t|/5$ for $t\ne0$, and is zero at $t=0$.
Together with $\phi(0)/6<1/15$ and the strict scale-remainder bound,
this proves the assertion for both signs and every threshold. In
particular there is no restriction to sampled thresholds or to a
bounded part of a support ray.
\end{proof}

\subsection{Finite bands and the exact support polynomial}\label{ms:legacy-support:polynomial}
For integers $36\le N\le L$ and tested radii $u\ge5/2$, use the
simultaneous endpoint bounds
\[
 B_+=\frac{u^4-6u^2+3}{24\sqrt{N-1}},\quad
 D_+=\frac{u^5-15u}{120(N-1)},\quad
 D_-=\frac{u^5-15u}{120(L-1)}.
\]
For an infinite index range, set $D_-=0$ instead. Put
\[
 W_*=A-15D_+/2,\qquad
 T_*=A-3D_-+\frac{9B_+^2}{4W_*}.
\]
The proof of Lemma~\ref{ms:legacy-support:envelope} applies uniformly on the whole band
if $W_*>0$ and $T_*\ge A/2+2B_++3D_+$. The restriction $u\ge5/2$
ensures that all the displayed $A,B_+,D_+,D_-$ are nonnegative. In particular the favorable
term is frozen at the upper index endpoint, and the adverse terms at
the lower one. No monotonicity of a signed coefficient is assumed.

Let $(P,H)$ be an attained positive closed maximizing pair, standardized
to mean zero and variance one. Write $C=C_n\ge \CE$,
$\rho=\mathbb E|X|^3$, $\beta=\mathbb E[X|X|]$, and
$d=n(C_{n-1}-C_n)_+$. The finite contact/predecessor bound is
\[
 \rho I(t)\le\mathcal H_n(H,t)+
 C\{-|t|^3+3\rho t^2/2+3\beta t+\rho e_n\}+\rho s_nd,
\]
with $I=0$ on the entire closed support. Suppose also, for a finite
band, that $C_n-\CE\ge K/n$, where $K>0$. The published finite cap gives
\[
 c=0.4095+K/L<C_n,\qquad
 \rho<r=\min\left\{\frac{0.3031(0.646)}{c-0.3031},
                        \frac{0.3322(0.429)}{c-0.3322}\right\}.
\]
For the unrestricted nonnegative case put $K=0$, $c=0.4095$.
The sign-moment inequality
$\rho\ge(1+\beta^2)/\sqrt{1-\beta^2}$ supplies $|\beta|<b$ whenever
$(1+b^2)^2>r^2(1-b^2)$ and $0<b<1$. All these statements concern the
same current law. The lower endpoint cannot replace $L$ in $K/L$.

When a preceding resource enclosure supplies sharper actual bounds
$\rho\le r$ and $|\beta|\le b$, these may instead be used in the
same polynomials. Their validity is the only extra premise.

Use positive rational $\ell,S$ with $\ell^2\le N-1$ and $S^2>N/(N-1)$. For $u=|t|$ define
\[
 Q(u)=-u^3+3ru^2/2+3bu+0.384r/N.
\]
First check $Q<0$, before replacing $C$ by its lower bound $c$.
Replace $1/\sqrt{N-1}$ by $1/\ell$ in the adverse remainder and use
$\widehat B_+=(u^4-6u^2+3)/(24\ell)\ge B_+$ in the linear term of
the domination gate. Explicitly, a sufficient rational remainder is
\[
 E_*(u)=\frac{51u^6}{14400(N-1)\ell}
       +\frac{51(3u^2-2)}{960(N-1)\ell}
       +\frac{u}{8(N-1)^2}.
\]
The finite-band formulas are used only where their denominators and
listed sign gates have been proved. Strict negative contact for every $n$ in the band follows
from the polynomial inequality
\[
 \frac25S^3\left\{(A-3D_-)W_*+
              \frac{(u^4-6u^2+3)^2}{256(N-1)}\right\}
 +S^3E_*W_*+\{cQ+rSq\}W_*<0.
\]
The second polynomial gate is $W_*>0$. The third is
\[
 \{A/2-3D_--3D_+-2\widehat B_+\}W_*+
         \frac{(u^4-6u^2+3)^2}{256(N-1)}>0.
\]
This is a rational sufficient version of the cleared-denominator
domination hypothesis; its quadratic $B_+^2$ term still uses the exact
factor $1/(N-1)$.
The bound on $T_*$ is positive, so the replacement $\phi_0<0.4$ and
$s_n<S$ both have their correct direction.

The retained exact acceptance condition tests these polynomials on all adjacent closed intervals
of length $0.05$ from $R$ to $4.5$. Positive Bernstein coefficients prove
their signs on every point of each cell, including endpoints. On the
rays $u\ge4.5$ use Lemma~\ref{ms:legacy-support:far-gaussian},
$|\mathcal H_n(H,t)|< u^3/15+u/5+1/2$.
The ray certificate must prove that the resulting cubic and its first three
derivatives are strictly negative at $4.5$, and that the coefficient $Q$
and its derivative gates permit replacing $C$ by $c$ on the whole ray.
This excludes the whole rays. Thus, if $d\le q$, actual support is in
$(-R,R)$, including exclusion of arbitrarily small exterior atoms and
exterior accumulation points.

\subsection{Retained support rows and their application domains}
\label{ms:legacy-support:rows}
The older local route specified $K=0$, $R=3.84$, $q=0.19$ on
\[
[45,49],\ [50,59],\ [60,89],\ [90,134],\ [135,199],\ [200,\infty).
\]
That route used only indices at least 47; this paragraph records its
literal support premises and makes no global exclusion claim. The last row uses $D_-=0$;
there is no upper index restriction. A complete acceptance must reconstruct all closed $u$ cells and both exterior rays. Missing an index band
or the infinite row is an acceptance failure.

The older finite-index route specified five further support rows:
\[
\begin{array}{c|c|c|c}
[N,L]&K&R&q\\\hline
{[40,40]}&1.5&3.3&0.88\\
{[41,49]}&1.5&2.9&1.5/40\\
{[45,45]}&0.7&3.5&0.8341\\
{[46,98]}&0.7&3.5&0.7/45\\
{[46,46]}&0.54&3.4&79/450
\end{array}
\]
The second and fourth are used only on $41\le n\le45$ and
$46\le n\le50$, respectively. Retaining their originally certified
upper endpoints makes their inequalities weaker, so restricting the
application introduces no new arithmetic assertion. These are support
implications, not yet exclusions: the separately proved and accepted
resource calculation would supply the latter. They are not an independent
proof of Assumption~\ref{ms:base41}.

\subsection{One bounded-threshold series}\label{ms:legacy-support:series}
Put $c_k=(-1)^k/[2^kk!(2k+1)]$, and let $(x)_j$ denote the rising
factorial. Define the following rational polynomials:
\begin{align*}
 P_j(H,t)={}&\sum_{k=0}^{j+1}
 \frac{c_k(k+1/2)_{j+1-k}}{(j+1-k)!}(H-t)^{2k+1}
       -c_{j+1}H^{2j+3}\\
 &+\frac{(-1)^{j+1}tH^{2j+2}}{2^{j+1}(j+1)!}
       +\frac{(-1)^jH^{2j+1}(t^2-1)}{2^{j+1}j!}.
\end{align*}
In particular,
\[
 P_0=(t^3-3t)/6,\qquad
 P_1=\frac{3t-t^3}{4}H^2+\frac{t^4-6t^2+3}{8}H
                  +\frac{-t^5+10t^3-15t}{40}.
\]
Fix $h>0$, $E(x)=(193/71)^{\lceil x\rceil}$, and
\[
 B_r=\frac{7(h+r)}6E((h+r)^2/6)
       +\left(h+r+\frac{h(r^2+1)}8\right)E(h^2/8).
\]
\begin{lemma}\label{ms:legacy-support:finite-series}
For every $n\ge4$, $|H|\le h$ and $|t|\le r$,
\[
 \left|\mathcal H_n/\phi(0)-\sum_{j=0}^5P_j/n^j\right|
       \le4^7B_r/n^6.
\]
\end{lemma}
\begin{proof}
Put $z=1/n$. Integration of the Gaussian exponential series gives
$\mathcal H_n/\phi(0)=S(z)/z$, where
\begin{align*}
 S(z)={}&\sum_{k\ge0}c_k\{(H-t)^{2k+1}(1-z)^{-k-1/2}
                                  -H^{2k+1}\}z^k\\
 &+te^{-H^2z/2}+\tfrac12zH(t^2-1)e^{-H^2z/2}.
\end{align*}
Its constant coefficient vanishes, and expansion of the binomial factor
gives $[z^{j+1}]S=P_j$. At $z_0=1/4$, the sum obtained by replacing
all coefficients by their absolute-value majorants is at most $B_r$:
use $(1-z_0)^{-1/2}<7/6$, $1/(2k+1)\le1$, and $e<193/71$.
The latter bound follows, for example, by summing the exponential series
through degree six and bounding the rest by a geometric series with
ratio $1/8$. All products and rearrangements are therefore absolutely
convergent. The discarded terms have degrees at least seven in $S$,
so their sum at $0\le z\le z_0$ is at most $(z/z_0)^7B_r$.
Divide by $z$ to get the result. This proves a finite uniform bound,
including every threshold in the stated interval and every $n\ge4$.
\end{proof}

\subsection{A common contact test for radius and strip refinement}
\label{ms:legacy-support:panels}
Suppose $n\ge47$, $0\le C-\CE\le0.7/n$,
$d=n(C_{n-1}-C)_+\le0.19$ and
$|X|\le R_0$. A previously acquired row gives
\[
 \rho\le r,\quad 0\le b_l\le\beta\le b_u<1,\quad |H|\le h.
\]
Define $Q(t)=-|t|^3+3\rho t^2/2+3\beta t$ and
$L_P(t)=\CE Q(t)+\phi(0)(t^3-3t)/6$.
For $t>0$, upper polynomials for $L_P(t),Q(t)$ are
\[
 (1/15-0.4095)t^3+1.5(0.40974)rt^2+(3(0.40974)b_u-0.3988/2)t,
 \quad -t^3+1.5rt^2+3b_ut.
\]
For $t<0$, write $u=|t|$ and use
\[
 (-0.3988/6-0.4095)u^3+1.5(0.40974)ru^2+(0.2-3(0.4095)b_l)u,
 \quad -u^3+1.5ru^2-3b_lu.
\]
These coefficient choices are valid because $b_l\ge0$ is checked first.
On any closed cell of one sign, let $L_J$ be an upper bound for the
first polynomial and let $Q_J\ge0$ bound the positive part of the second.
For $j=1,2,3$, take $m_j$ to be the positive part of the largest tensor
Bernstein coefficient of $P_j$ on $[-h,h]\times J$. For $j=4,5$, the
absolute monomial bound
$\sum_{a,b}|[H^at^b]P_j|h^aR_J^b$ suffices, where
$R_J=\max_{t\in J}|t|$. The support contact is then bounded above by
\begin{equation}\label{ms:legacy-support:paneltest}
 L_J+\frac{0.7Q_J}{47}
 +0.4\left(\sum_{j=1}^5\frac{m_j}{47^j}
                +\frac{4^7B_{R_J}}{47^6}\right)
 +r\left(0.19S+\frac{0.384(0.425)}{47}\right),
 \qquad S>\sqrt{47/46}.
\end{equation}
Indeed substitute the finite series in contact, split $C=\CE+(C-\CE)$,
and use $C<0.40974+0.7/47<0.425$. Every error retained in
\eqref{ms:legacy-support:paneltest} is nonnegative and decreases with $n$.
A strictly negative rational result therefore excludes the entire cell
for every integer $n\ge47$, including atoms and accumulation points.

For a coordinate strip, conditional centering gives
\[
 X=\ell(Y-p),\quad
 \ell=\sqrt{(1-\theta)/(p(1-p))},\quad
 \beta=(1-\theta)(1-2p)+\kappa,\quad |\kappa|\le\theta.
\]
If $p\in[p_l,p_u]\subset(0,1/2)$ and $\theta\le\Theta$, put
$\ell_l\le\sqrt{(1-\Theta)/(p_u(1-p_u))}$.
The forbidden strip $Y\notin(-0.55,1.55)$ maps into
\[
 X\le-\ell_l(p_l+0.55),\qquad X\ge\ell_l(1.55-p_u).
\]
Apply \eqref{ms:legacy-support:paneltest} between these thresholds and $\pm R_0$.
An interval beyond the known radius is empty; an equal endpoint is a
singleton that is still checked. If only an upper $\beta$ bound is given,
use $b_l=(1-\Theta)(1-2p_u)-\Theta$, after checking its nonnegativity.

The retained schedule lists the following conditional consequences; each
requires the complete numerical tests just described:
\[
\begin{array}{l|c|c|r}
\text{input rows}&\text{prior radius}&\text{conclusion}&\text{closed cells}\\\hline
10\text{ central preparation rows}&3.84&|X|<2.5&1080\\
45\text{ central rows}&2.5&-0.55<Y<1.55&1431\\
38\text{ left rows}&2.7&-0.55<Y<1.55&3073
\end{array}
\]
The corresponding upper contact bounds must be strictly negative on
every listed cell. Each input moment interval must first be established
by the preceding resource estimate.
The left radii are applications of the all-threshold support polynomial,
with $(N,L,K,q)=(47,\infty,0,0.19)$ and the actual moment caps
\[
 (r,b,R)=(1.390,0.430,3),\qquad (1.2563,0.3926,2.7).
\]
Conditional on the corresponding polynomial checks, the same heat proof
supplies the initial radius, both left contractions, and the five finite-base rows. The bounded-threshold
series supplies the central contraction and both coordinate strips.

\subsection{Relation to the zero-loss bound}
\label{ms:legacy-support:retention}

The $65$ contractions of Section~\ref{ms:support} use zero predecessor
loss and moment bounds established in the stated order. The index-$27$
condition, the five positive-excess finite-band cases and their support
strips have different hypotheses. Their proofs require the separate
polynomial inequalities and closed-cell bounds given above.
